%%**************************************************************
%% Vorlage fuer Bachelorarbeiten (o.ä.) der DHBW
%%
%% Autor: Tobias Dreher, Yves Fischer, Michael Gruben, Markus Barthel
%% Datum: 06.07.2011 - 22.08.2014
%%
%% Autor: Ferdinand König und Maximilian Köhler
%% Datum: 2020 - 2022
%%**************************************************************

%!TEX root = ../main.tex

%
% Nahezu alle Einstellungen koennen hier getaetigt werden
%

\RequirePackage[l2tabu, orthodox]{nag}	% weist in Commandozeile bzw. log auf veraltete LaTeX Syntax hin

\documentclass[%
    % final,
	pdftex,
	oneside,			% Einseitiger Druck (oneside) oder zweiseitig (twoside)
	% headings=openright,	% Kapitelanfänge immer auf rechter Seite (bei zweiseitig)
	cleardoublepage=empty,	% Leere Vakatseiten
	12pt,				% Schriftgroesse
	parskip=half,		% Halbe Zeile Abstand zwischen Absätzen (half).
%	topmargin = 10pt,	% Abstand Seitenrand (Std:1in) zu Kopfzeile [laut log: unused]
	headheight = 20pt,	% Höhe der Kopfzeile
%	headsep = 30pt,	% Abstand zwischen Kopfzeile und Text Body  [laut log: unused]
	headsepline,		% Linie nach Kopfzeile.
	% footsepline,		% Linie vor Fusszeile.
	footheight = 20pt,	% Höhe der Fusszeile
	abstracton,		% Abstract Überschriften
	DIV=calc,		% Satzspiegel berechnen
	headinclude=false,	% Kopfzeile nicht in den Satzspiegel einbeziehen
	footinclude=false,	% Fußzeile nicht in den Satzspiegel einbeziehen
	listof=totoc,		% Abbildungs-/ Tabellenverzeichnis im Inhaltsverzeichnis darstellen
	toc=bibliography,	% Literaturverzeichnis im Inhaltsverzeichnis darstellen
	pointlessnumbers,
	% chapterprefix=true,
	% bibliography=openstyle
]{scrreprt}	% Koma-Script report-Klasse (scrreprt), fuer laengere Bachelorarbeiten alternativ auch: scrbook

\raggedbottom

% Einstellungen laden
\usepackage{xstring}
\usepackage[utf8]{inputenc}
\usepackage[T1]{fontenc}

\newcommand{\einstellung}[1]{%
  \expandafter\newcommand\csname #1\endcsname{}
  \expandafter\newcommand\csname setze#1\endcsname[1]{\expandafter\renewcommand\csname#1\endcsname{##1}}
}
\newcommand{\langstr}[1]{\einstellung{lang#1}}

\input{ads/einstellungen_liste} % verfügbare Einstellungen
%%%%%%%%%%%%%%%%%%%%%%%%%%%%%%%%%%%%%%%%%%%%%%%%%%%%%%%%%%%%%%%%%%%%%%%%%%%%%%%
%                                   Einstellungen
%
% Hier können alle relevanten Einstellungen für diese Arbeit gesetzt werden.
% Dazu gehören Angaben u.a. über den Autor sowie Formatierungen.
%%%%%%%%%%%%%%%%%%%%%%%%%%%%%%%%%%%%%%%%%%%%%%%%%%%%%%%%%%%%%%%%%%%%%%%%%%%%%%%


%%%%%%%%%%%%%%%%%%%%%%%%%%%%%%%%%%%% Sprache %%%%%%%%%%%%%%%%%%%%%%%%%%%%%%%%%%%
%% Aktuell sind Deutsch und Englisch unterstützt.
%% Es werden nicht nur alle vom Dokument erzeugten Texte in
%% der entsprechenden Sprache angezeigt, sondern auch weitere
%% Aspekte angepasst, wie z.B. die Anführungszeichen und
%% Datumsformate.
\setzesprache{de} % de oder en
%%%%%%%%%%%%%%%%%%%%%%%%%%%%%%%%%%%%%%%%%%%%%%%%%%%%%%%%%%%%%%%%%%%%%%%%%%%%%%%%



%%%%%%%%%%%%%%%%%%%%%%%%%%%%%%%%%%% Angaben  %%%%%%%%%%%%%%%%%%%%%%%%%%%%%%%%%%%
%% Die meisten der folgenden Daten werden auf dem
%% Deckblatt angezeigt, einige auch im weiteren Verlauf
%% des Dokuments.
\setzematrikelnr{23176975}
\setzekurs{TMB19 AM2}
\setzetitel{Critical clearing time of synchronous generators}
\setzedatumAbgabe{\today}
\setzefirma{Schaeffler Technologies AG \& Co. KG}
\setzefirmenort{Industriestraße 1-3, 91074 Herzogenaurach}
\setzeabgabeort{Erlangen}
\setzestudiengang{Energy Technology}
\setzedhbw{Mannheim}
\setzebetreuer{Dipl.-Phys. Philipp Hörning}
\setzegutachter{Prof. Dr.-Ing. Simon Thiele}
\setzezeitraum{27.06.2022\ -\ 19.09.2022}
\setzearbeit{Seminararbeit}
\setzeautor{Maximilian Köhler} 




%%%%%%%%%%%%%%%%%%%%%%%%%%%% Literaturverzeichnis %%%%%%%%%%%%%%%%%%%%%%%%%%%%%%
%% Bei Fehlern während der Verarbeitung bitte in ads/header.tex bei der
%% Einbindung des Pakets biblatex (ungefähr ab Zeile 110,
%% einmal für jede Sprache), biber in bibtex ändern.
\newcommand{\ladeliteratur}{%
\addbibresource{literatur.bib}
%\addbibresource{weitereDatei.bib}
}
%% Zitierstil
%% siehe: http://ctan.mirrorcatalogs.com/macros/latex/contrib/biblatex/doc/biblatex.pdf (3.3.1 Citation Styles)
%% mögliche Werte z.B numeric-comp, alphabetic, authoryear, numeric
\setzezitierstil{ieee}
%\setzezitierstil{alphabetic}
%\setzezitierstil{authoryear}




%%%%%%%%%%%%%%%%%%%%%%%%%%%%%%%%% Layout %%%%%%%%%%%%%%%%%%%%%%%%%%%%%%%%%%%%%%%
%% Verschiedene Schriftarten
% laut nag Warnung: palatino obsolete, use mathpazo, helvet (option scaled=.95), courier instead
\setzeschriftart{charter} % palatino oder goudysans, lmodern, libertine, mathptmx
% palatino oder lmodern ganz nett

%% Paket um Textteile drehen zu können
%\usepackage{rotating}
%% Paket um Seite im Querformat anzuzeigen
%\usepackage{lscape}

%% Seitenränder
% \setzeseitenrand{2.5cm}

%% Abstand vor Kapitelüberschriften zum oberen Seitenrand
\setzekapitelabstand{0pt}

%% Spaltenabstand
\setzespaltenabstand{10pt}
%%Zeilenabstand innerhalb einer Tabelle
\setzezeilenabstand{1.5}




%%%%%%%%%%%%%%%%%%%%%%%%%%%%% Verschiedenes  %%%%%%%%%%%%%%%%%%%%%%%%%%%%%%%%%%%
%% Farben (Angabe in HTML-Notation mit großen Buchstaben)
\newcommand{\ladefarben}{%
	\definecolor{LinkColor}{HTML}{00007A}
	\definecolor{ListingBackground}{HTML}{FCF7DE}
}

%%Mathematikpakete benutzen (Pakete aktivieren)
\usepackage{amsmath}
\usepackage{amssymb}

%% Programmiersprachen Highlighting (Listings)
\newcommand{\listingsettings}{%
	\lstset{%
		language=Java,			% Standardsprache des Quellcodes
		numbers=left,			% Zeilennummern links
		stepnumber=1,			% Jede Zeile nummerieren.
		numbersep=5pt,			% 5pt Abstand zum Quellcode
		numberstyle=\tiny,		% Zeichengrösse 'tiny' für die Nummern.
		breaklines=true,		% Zeilen umbrechen wenn notwendig.
		breakautoindent=true,	% Nach dem Zeilenumbruch Zeile einrücken.
		postbreak=\space,		% Bei Leerzeichen umbrechen.
		tabsize=2,				% Tabulatorgrösse 2
		basicstyle=\ttfamily\footnotesize, % Nichtproportionale Schrift, klein für den Quellcode
		showspaces=false,		% Leerzeichen nicht anzeigen.
		showstringspaces=false,	% Leerzeichen auch in Strings ('') nicht anzeigen.
		extendedchars=true,		% Alle Zeichen vom Latin1 Zeichensatz anzeigen.
		captionpos=b,			% sets the caption-position to bottom
		backgroundcolor=\color{ListingBackground}, % Hintergrundfarbe des Quellcodes setzen.
		xleftmargin=0pt,		% Rand links
		xrightmargin=0pt,		% Rand rechts
		frame=single,			% Rahmen an
		frameround=ffff,
		rulecolor=\color{darkgray},	% Rahmenfarbe
		fillcolor=\color{ListingBackground},
		keywordstyle=\color[rgb]{0.133,0.133,0.6},
		commentstyle=\color[rgb]{0.133,0.545,0.133},
		%stringstyle=\color[rgb]{0.627,0.126,0.941}
		stringstyle=\color{red}
	}
}




%%%%%%%%%%%%%%%%%%%%%%%%%%%%%%%% Eigenes %%%%%%%%%%%%%%%%%%%%%%%%%%%%%%%%%%%%%%%
%% Hier können Ergänzungen zur Präambel vorgenommen werden (eigene Pakete, Einstellungen)

%colorpackage
% \usepackage{color}
\usepackage[table,dvipsnames]{xcolor}


%listing package for defining a new language
\usepackage{listings}

% starting footnotes at 1
% \usepackage{perpage}
% \MakePerPage[1]{footnote}

%Functions
\def\bsq#1{%both single quotes
\lq{#1}\rq}


%%%%%%%%%%%%%%%%%%%%%%%%%%%%%%%%%%%%%%%%%%%%
% selfmade language pakets
% JavaScript language
\definecolor{lightgray}{RGB}{227, 227, 227}
\definecolor{darkgray}{RGB}{100, 100, 100}
\definecolor{purple}{rgb}{0.65, 0.12, 0.82}
\definecolor{schaeffler}{RGB}{0, 137, 61}
\definecolor{dark-green}{RGB}{0, 110, 93}
\definecolor{middle-green}{RGB}{115, 161, 149}
\definecolor{light-green}{RGB}{199, 222, 160}

\definecolor{pie1}{RGB}{115, 161, 149}
\definecolor{pie2}{RGB}{192, 198, 191}
\definecolor{pie3}{RGB}{135, 135, 135}
\definecolor{pie4}{RGB}{29, 155, 178}
\definecolor{pie5}{RGB}{182, 186, 194}
\definecolor{pie6}{RGB}{161, 200, 97}
\definecolor{pie7}{RGB}{67, 99, 91}
\definecolor{pie8}{RGB}{112, 123, 110}
\definecolor{pie9}{RGB}{113, 113, 113}

\lstdefinelanguage{JS}{
  keywords={break, case, catch, continue, debugger, default, delete, do, else, false, finally, for, function, if, in, instanceof, new, null, return, switch, this, throw, true, try, typeof, var, void, while, with},
  morecomment=[l]{//},
  morecomment=[s]{/*}{*/},
  morestring=[b]',
  morestring=[b]",
  ndkeywords={class, export, boolean, throw, implements, import, this},
  keywordstyle=\color{blue}\bfseries,
  ndkeywordstyle=\color{darkgray}\bfseries,
  identifierstyle=\color{black},
  commentstyle=\color[rgb]{0.133,0.545,0.133},
  %commentstyle=\color{purple}\ttfamily,
  stringstyle=\color{red}\ttfamily,
  sensitive=true
}
 % lese Einstellungen

\input{lang/strings} % verfügbare Strings
\input{lang/\sprache} % Übersetzung einlesen

% Einstellung der Sprache des Paketes Babel und der Verzeichnisüberschriften
\iflang{de}{\usepackage[english, ngerman]{babel}}
\iflang{en}{\usepackage[ngerman, english]{babel}}

%%%%%%% Package Includes %%%%%%%

\usepackage[left=3cm,right=2cm,top=2.5cm,bottom=2.5cm,foot=1cm]{geometry}	% Seitenränder und Abstände
\usepackage[activate]{microtype} %Zeilenumbruch und mehr
\usepackage[onehalfspacing]{setspace}
\usepackage{makeidx}
\usepackage[autostyle=true,german=quotes]{csquotes}
\usepackage{longtable}
\usepackage{enumitem}	% mehr Optionen bei Aufzählungen
\usepackage{graphicx}
\usepackage{pdfpages}   % zum Einbinden von PDFs
% \usepackage[table]{xcolor} 	% für HTML-Notation
\usepackage{float}
\usepackage{array}
\usepackage{calc}		% zum Rechnen (Bildtabelle in Deckblatt)
\usepackage[right]{eurosym}
\DeclareUnicodeCharacter{20AC}{\euro}
\usepackage{wrapfig}
\usepackage{pgffor} % für automatische Kapiteldateieinbindung
\usepackage[hang, multiple, stable]{footmisc} % Fussnoten; perpage für jede Seite
\usepackage{chngcntr}
\counterwithout{footnote}{chapter}
\usepackage[printonlyused]{acronym} % falls gewünscht kann die Option footnote eingefügt werden, dann wird die Erklärung nicht inline sondern in einer Fußnote dargestellt
\usepackage{listings}
\usepackage{tabularx}
% \usepackage{pdflscape}
\usepackage{lscape}
\usepackage{rotating}
\usepackage[labelfont={bf},font={small},format={hang},singlelinecheck={false}]{caption} % Einstellungen für Bildunterschriften/Tabellenüberschriften
\setcapwidth[c]{.9\textwidth}
\usepackage{subcaption} 
\usepackage{ltxtable}
% \usepackage{filecontents}
\setlength{\skip\footins}{10pt plus 6pt minus 0pt} % Abstand zwischen Fußnoten und Fließtext erhöhen. Latex-Standard: 10pt plus 4pt minus 2pt
\usepackage{mathtools}
\usepackage{tikz}
\usetikzlibrary{positioning}
\usepackage{pgf-pie} % https://www.namsu.de/Extra/pakete/Pie_Chart.html
\usepackage{pgfplots}
\pgfplotsset{compat=1.16}
\usepackage{multirow}
\usepackage[absolute]{textpos}
\usepackage{tcolorbox}
% \usepackage{minitoc}
\usepackage{afterpage}
% \usepackage[Glenn]{fncychap}
% Sonny, Lenny, Glenn, Conny, Rejne, Bjarne, Bjornstrup

% Notizen. Einsatz mit \todo{Notiz} oder \todo[inline]{Notiz}.
\usepackage[obeyFinal,backgroundcolor=yellow,linecolor=black]{todonotes}
% Alle Notizen ausblenden mit der Option "final" in \documentclass[...] oder durch das auskommentieren folgender Zeile
% \usepackage[disable]{todonotes}

% Kommentarumgebung. Einsatz mit \comment{}. Alle Kommentare ausblenden mit dem Auskommentieren der folgenden und dem aktivieren der nächsten Zeile.
% \newcommand{\comment}[1]{\par {\bfseries \itshape \color{blue} #1 \par}} %Kommentar anzeigen
%\newcommand{\comment}[1]{} %Kommentar ausblenden


%%%%%% Configuration %%%%%

%% Anwenden der Einstellungen

\usepackage{\schriftart}
% Überschriften auch in gesetzter Schriftart
\setkomafont{disposition}{%
	\normalfont\bfseries
}
\setkomafont{dictum}{\normalfont}
% Verwendung der Schrift ohne Serifen
% \renewcommand*{\familydefault}{\sfdefault}
% \addtokomafont{disposition}{\sffamily}

\ladefarben{}

% Titel, Autor und Datum
\title{\titel}
\author{\autor}
\date{\datum}

% PDF Einstellungen
\usepackage[%
	pdftitle={\titel},
	pdfauthor={\autor},
	pdfsubject={\arbeit},
	pdfcreator={pdflatex, LaTeX with KOMA-Script},
	pdfpagemode=UseOutlines, 		% Beim Oeffnen Inhaltsverzeichnis anzeigen
	pdfdisplaydoctitle=true, 		% Dokumenttitel statt Dateiname anzeigen.
	pdflang={\sprache}, 			% Sprache des Dokuments.
	%hidelinks,						% entfernt Umrandung von verlinkten Stellen, ohne Verlinkung zu löschen
]{hyperref}

% (Farb-)einstellungen für die Links im PDF
\hypersetup{%
	colorlinks=true, 		% Aktivieren von farbigen Links im Dokument
	linkcolor=black, 	% Farbe festlegen
	citecolor=black,
	filecolor=black,
	menucolor=black,
	urlcolor=black,
	linktocpage=true, 		% Nicht der Text sondern die Seitenzahlen in Verzeichnissen klickbar
	bookmarksnumbered=true 	% Überschriftsnummerierung im PDF Inhalt anzeigen.
}
% Workaround um Fehler in Hyperref, muss hier stehen bleiben
\usepackage{bookmark} %nur ein latex-Durchlauf für die Aktualisierung von Verzeichnissen nötig

% Schriftart in Captions etwas kleiner
\addtokomafont{caption}{\small}

% Literaturverweise (sowohl deutsch als auch englisch)
\iflang{de}{%
\usepackage[
	backend=biber,		% empfohlen. Falls biber Probleme macht: bibtex
	bibwarn=true,
	bibencoding=utf8,	% wenn .bib in utf8, sonst ascii
	% sortlocale=de_DE,
	sorting=none,		% Altenativen: https://tex.stackexchange.com/questions/51434/biblatex-citation-order
	style=\zitierstil,
]{biblatex}
}
\iflang{en}{%
\usepackage[
	backend=biber,		% empfohlen. Falls biber Probleme macht: bibtex
	bibwarn=true,
	bibencoding=utf8,	% wenn .bib in utf8, sonst ascii
	% sortlocale=en_US,
	sorting=none,
	style=\zitierstil,
]{biblatex}
}

\setcounter{biburlnumpenalty}{100}
\setcounter{biburlucpenalty}{100}
\setcounter{biburllcpenalty}{100}

\ladeliteratur{}

% Glossar
\usepackage[nonumberlist,toc,automake]{glossaries}
\addtokomafont{descriptionlabel}{\normalfont\bfseries}

%Kopf- und Fußzeilen
\usepackage[plainfootsepline=yes]{scrlayer-scrpage}

%%%%%% Additional settings %%%%%%

% Hurenkinder und Schusterjungen verhindern
% http://projekte.dante.de/DanteFAQ/Silbentrennung
\clubpenalty = 10000 % schließt Schusterjungen aus (Seitenumbruch nach der ersten Zeile eines neuen Absatzes)
\widowpenalty = 10000 % schließt Hurenkinder aus (die letzte Zeile eines Absatzes steht auf einer neuen Seite)
\displaywidowpenalty=10000

% Bildpfad
\graphicspath{{images/}}

% Einige häufig verwendete Sprachen
\lstloadlanguages{PHP,Python,Java,C,C++,bash}
\listingsettings{}
% Umbennung des Listings
\renewcommand\lstlistingname{\langlistingname}
\renewcommand\lstlistlistingname{\langlistlistingname}
\def\lstlistingautorefname{\langlistingautorefname}

% Abstände in Tabellen
\setlength{\tabcolsep}{\spaltenabstand}
\renewcommand{\arraystretch}{\zeilenabstand}

% Anhangsverzeichnis
\makeatletter% --> De-TeX-FAQ
% Weitergabe des folgenden Codes oder Modifikationen davon nur unter Nennung
% der Originalquelle: <http://www.komascript.de/comment/1073#comment-1073>,
% gestattet.
% Leistungsfähigere Lösung unter <https://komascript.de/comment/5578#comment-5578>.

% Inhaltsverzeichnis für den Anhang erstellen 
% \newcommand*{\maintoc}{% Hauptinhaltsverzeichnis
%   \begingroup
%     \@fileswfalse% kein neues Verzeichnis öffnen
%     \renewcommand*{\appendixattoc}{% Trennanweisung im Inhaltsverzeichnis
%       \value{tocdepth}=-10000 % lokal tocdepth auf sehr kleinen Wert setzen
%     }%
%     \tableofcontents% Verzeichnis ausgeben
%   \endgroup
% }
% \newcommand*{\appendixtoc}{% Anhangsinhaltsverzeichnis
%   \begingroup
%     \edef\@alltocdepth{\the\value{tocdepth}}% tocdepth merken
%     \setcounter{tocdepth}{-10000}% Keine Verzeichniseinträge
%     \renewcommand*{\contentsname}{% Verzeichnisname ändern
%       \langanhang}%
%     \renewcommand*{\appendixattoc}{% Trennanweisung im Inhaltsverzeichnis
%       \setcounter{tocdepth}{\@alltocdepth}% tocdepth wiederherstellen
%     }%
%     \tableofcontents% Verzeichnis ausgeben
%     \setcounter{tocdepth}{\@alltocdepth}% tocdepth wiederherstellen
%   \endgroup
% }
% \newcommand*{\appendixattoc}{% Trennanweisung im Inhaltsverzeichnis
% }
% \g@addto@macro\appendix{% \appendix erweitern
%   \if@openright\cleardoublepage\else\clearpage\fi% Neue Seite
%   \phantomsection
%   \addcontentsline{toc}{chapter}{\appendixname}% Eintrag ins Hauptverzeichnis
%   \addtocontents{toc}{\protect\appendixattoc}% Trennanweisung in die toc-Datei
% }
% \makeatother

% Kreisdiagramme
\def\printonlylargeenough#1#2{\unless\ifdim#2pt<#1pt\relax
#2\printnumbertrue
\else
\printnumberfalse
\fi}
\newif\ifprintnumber

\newcommand{\tab}{~~~~~}

\clearpairofpagestyles
\ohead[]{\textsc{\headmark}}				% Kopfzeile außen immer mit Headmark versehen
\automark[section]{chapter}		% Headmark bestehend aus Kolumnentitel
\ofoot[\pagemark]{\pagemark}	% Fußzeile mit Seitenzahl außen
\renewcommand*\chapterpagestyle{plain.scrheadings}
% \renewcommand*\partpagestyle{plain.scrheadings}		%Bei Verwendung von Parts als Überschriftenebene: Setzen des Pagestyles global



% \makeglossaries

\begin{document}
	\pagenumbering{Roman}

	% Deckblatt
	\begin{spacing}{1}
		\begin{titlepage}
	\thispagestyle{empty}
	\setcounter{page}{-1}
	
	% \begin{textblock*}{\textwidth}(30mm,33mm)
	% \includegraphics[width=47mm]{ees}
	% \end{textblock*}
	
	\begin{textblock*}{\textwidth}(127mm,192mm)
	\includegraphics[width=93mm]{fausiegel}
	\end{textblock*}
	
	% {\raggedleft
	% \textsc{Lehrstuhl für Elektrische Energiesysteme}\\
	% Vorstand: Univ.-Prof. Dr.-Ing. Matthias Luther
	% \par}

	\begin{minipage}[]{47mm}
		\includegraphics[width=47mm]{ees}
	\end{minipage}
	\begin{minipage}[]{.7\textwidth}
		\raggedleft
		\textsc{Lehrstuhl für Elektrische Energiesysteme}\\
		Vorstand: Univ.-Prof. Dr.-Ing. Matthias Luther
	\end{minipage}
	
	\vspace{51mm}
	
	{\centering
	
	%##########################################################
	% BEARBEITUNGSFELDER
	\large{Seminararbeit} \\	
	\Large{Critical clearing time of synchronous generators}
	%##########################################################
	
	\par}
	
	\vspace{105mm}
	
	{\raggedright
	\begin{tabbing}
	XX \= XXXXXXXXXXXXXX \= XXXXXXXXXXXXXXXXXXXXXX \kill
	%##########################################################
	% BEARBEITUNGSFELDER
			\> \textbf{Bearbeiter:} 	\> B. Eng. Maximilian Köhler \\
			\>						\> 23176975		\\
	 		\>												\>									\\
			\> \textbf{Betreuer:}		\> M. Sc. Ilya Burlakin \\
			\>												\>									\\
			\> \textbf{Abgabedatum:}	\> \today															
	%##########################################################
	\end{tabbing}
	\par}
	% \cleardoublepage
\end{titlepage}

	\end{spacing}
	% \newpage
	\newpage
	% \cleardoublepage

	% stellt Abstand vor Kapitelüberschriften ein
	\RedeclareSectionCommand[beforeskip=\kapitelabstand         ]{chapter}
	% \newgeometry{left=3cm,right=2cm,top=2.5cm,bottom=2.5cm}

    \pagestyle{scrheadings}

    \setcounter{page}{2}
	% Sperrvermerk
	% %!TEX root = ../main.tex

% \newgeometry{left=2.5cm,right=2.5cm,top=2.5cm,bottom=2.5cm}

% Sperrvermerk direkt hinter Titelseite

\addchap*{\langsperrvermerk}
\thispagestyle{empty}

\vfill

\begin{figure}[H]
    \centering
    \href{https://www.kununu.com/de/schaeffler}{\includegraphics[width=12cm]{images/essential/schaeffler.png}}
\end{figure}
\vspace*{3cm}

\iflang{de}{
Der Inhalt dieser Bachelorarbeit darf weder als Ganzes noch in Auszügen Personen außerhalb des Prüfungsprozesses und des Evaluationsverfahrens zugänglich gemacht werden, sofern keine anders lautende Genehmigung des Unternehmens Schaeffler vorliegt.\\

\vspace*{1cm}
Insbesondere ist eine Weitergabe an Wettbewerber von Schaeffler oder eine Veröffentlichung, auch auszugsweise, nicht gestattet.\\

\vspace*{1cm}
Ausnahmen bedürfen der schriftlichen Genehmigung durch den entsprechenden Hauptabteilungsleiter von Schaeffler.
}

\iflang{en}{%
  The {\arbeit} on hand 
  \begin{center}{\itshape{} \titel{}\/}\end{center} 
   contains internal resp.\ confidential data of {\firma}. It is intended solely for inspection by the assigned examiner, the head of the {\studiengang} department and, if necessary, the Audit Committee \langanderdh{} {\dhbw}. It is strictly forbidden
    \begin{itemize}
    \item to distribute the content of this paper (including data, figures, tables, charts etc.) as a whole or in extracts,
    \item to make copies or transcripts of this paper or of parts of it,
    \item to display this paper or make it available in digital, electronic or virtual form.
    \end{itemize}
  Exceptional cases may be considered through permission granted in written form by the author and {\firma}.
}

\vspace{4cm}

% \restoregeometry
    % \newpage
	
	% Erklärung
 	%!TEX root = ../main.tex

\addchap*{\langerklaerung}
% \thispagestyle{empty}

\vspace*{1.5cm}

\begin{center}
    \begin{tabular}{| p{0.95\textwidth} |}
        \hline
        I certify that I have prepared this \arbeit without outside help and without using sources other than those specified and that the thesis has not been submitted in the same or a similar form to any other examination authority and has not been accepted by them as part of an examination. All statements that have been copied verbatim or in spirit are marked as such.\\
        \vspace{.5cm}
        Erlangen, \today \\ % \datumAbgabe\\
        \vspace*{.5cm}
        \singlespacing
        \rule{7cm}{.5pt}\\
        \autor\\[12pt]
        \hline
    \end{tabular}
\end{center}

\vfill

\begin{flushright}
    \begin{minipage}[]{0.7\textwidth}
        \textbf{Note:}\\[6pt]
        For reasons of readability, the generic masculine is primarily used in this \arbeit. Female and other gender identities are explicitly included where this is necessary for the statement.
    \end{minipage}
\end{flushright}

\vspace{2cm}
	%!TEX root = ../main.tex

\chapter*{Assignment of the paper}
%\vspace*{4cm}
\thispagestyle{plain.scrheadings}
{\large \textbf{Topic:} \parbox[t]{0.8\textwidth}{\titel}}
\newline

The \ac{CCT} is an essential parameter in power system stability
analysis. For example, in the case of \ac{SG}, the \acs{CCT} determines the
maximum fault-clearing time a generator can withstand without losing synchronism.
This seminar will introduce the concept of \acs{CCT} computing. We will discuss the factors
influencing \acs{CCT}, such as generator characteristics, system parameters, and fault type,
and explore the methods used to calculate \acs{CCT} in practical power system analysis.

% \begin{itemize}
%     \item Swing equation of synchronous generators
%     \item Lösen der Swing equation mithilfe von Python -> Lösung einer DGL zweiter Ordnung
%     \item Equal-area criterion -> Herleitung der Gleichungen
%     \item Simulation eines Fehlers -> Anwendung des equal-area criterions mithilfe von python.
%     \item Vergleich zwischen analytischer und simulativer Ergebnisse
% \end{itemize}

The seminar research paper should contain:
\begin{itemize}
    \item A literature research of governing equations describing the short-term dynamic behavior of \acs{SG}, relevant fault types and their influence;
    \item an investigation of the influences from machine characteristics and system parameters on the \acs{CCT};
    \item a computed simulation model for numerical determination of the \acs{CCT} with the \ac{EAC};
    \item simulations of system faults and comparisons to analytical solutions.
\end{itemize} 

	% Inhaltsverzeichnis
	\begin{spacing}{1.2}
		\begingroup
			% auskommentieren für Seitenzahlen unter Inhaltsverzeichnis
			\renewcommand*{\chapterpagestyle}{empty}
			\pagestyle{empty}

			%\setcounter{tocdepth}{1}
			%für die Anzeige von Unterkapiteln im Inhaltsverzeichnis
			\setcounter{tocdepth}{2}

			% \tableofcontents
			\maintoc
			\clearpage
		\endgroup
	\end{spacing}
	\newpage
    % \adjustmtc
    \pagestyle{scrheadings}
	
	% noch ausstehende ToDo's - Liste
	% \listoftodos

	% Abbildungsverzeichnis
	% \clearpage
	% \listoffigures

	% Tabellenverzeichnis
	% \clearpage
	% \listoftables
	
	% \cleardoublepage
	\clearpage
	\pagenumbering{arabic}
    
    \pagestyle{scrheadings}		% Kopf- und Fußzeile wie zuvor eingestellt
	
	%\setcounter{footnote}{1}	% Dont start at 2

	% Inhalt
	\foreach \i in {01,02,03,04,05,06,07,08,09,...,99} {%
		\edef\FileName{content/\i kapitel}%
			\IfFileExists{\FileName}{%
				\input{\FileName}
			}
			{%
				%file does not exist
			}
	}

	\clearpage
	
	\pagenumbering{Roman}
    \setcounter{page}{7}

	% Abkürzungsverzeichnis
	\clearpage
	%!TEX root = ../main.tex

%%%%%%%%%%%%%%%%%%%%%%%%%%%%%%%%%%%%%%%%%%%%%%%%%%%%%%%%%%%%%%%%
% Anmerkungen zur Verwendung:
%%%%%%%%%%%%%%%%%%%%%%%%%%%%%%%%%%%%%%%%%%%%%%%%%%%%%%%%%%%%%%%%
%
% nur verwendete Akronyme werden letztlich im Abkürzungsverzeichnis des Dokuments angezeigt
% Verwendung: 
%		\ac{Abk.}   --> fügt die Abkürzung ein, beim ersten Aufruf wird zusätzlich automatisch die ausgeschriebene Version davor eingefügt bzw. in einer Fußnote (hierfür muss in header.tex \usepackage[printonlyused,footnote]{acronym} stehen) dargestellt
%		\acs{Abk.}   -->  fügt die Abkürzung ein
%		\acf{Abk.}   --> fügt die Abkürzung UND die Erklärung ein
%		\acl{Abk.}   --> fügt nur die Erklärung ein
%		\acp{Abk.}  --> gibt Plural aus (angefügtes 's'); das zusätzliche 'p' funktioniert auch bei obigen Befehlen
%	siehe auch: http://golatex.de/wiki/%5Cacronym
%
%%%%%%%%%%%%%%%%%%%%%%%%%%%%%%%%%%%%%%%%%%%%%%%%%%%%%%%%%%%%%%%%

\addchap{\langabkverz}

\begin{acronym}[mmmmmm] %hier längstes Acro
% \begin{doublespacing}
% \setlength{\itemsep}{-\parsep}

\acro{ASR}{Area Specific Resistance}
\acro{CCM}{Catalyst Coated Membrane}
\acrodefplural{CCM}[CCMs]{Catalyst Coated Membranes}
\acro{ECD}{Electrochemical Deposition}
\acro{ESS}{energy storage system}
\acrodefplural{ESS}[ESS]{energy storage systems}
\acro{MEA}{Membrane Electrode Assembly}
\acrodefplural{MEA}[MEAs]{Membrane Electrode Assemblys}
\acro{PEM}{proton exchange membrane}
\acro{PEMWE}{proton exchange membrane water electrolysis}
\acro{PTL}{porous transport layer}
\acro{PVD}{Physical Vapour Deposition}
\acro{SOC}{State of Charge}

\end{acronym}

	% Glossar
	% \clearpage
	% %!TEX root = ../main.tex

%
% vorher in Konsole folgendes aufrufen:
%	makeglossaries makeglossaries dokumentation.acn && makeglossaries dokumentation.glo
%
%
% Glossareintraege --> referenz, name, beschreibung
% Aufruf mit \gls{...}
%
%\newglossaryentry{}{name={},plural={},description={}}
\newglossaryentry{FIP}{name={Forschungs- und Innovationsprojekt},description={Entwicklungsprojektart in der Schaeffler Gruppe mit dem Fokus, Projekte im Innovationsbereich möglichst frei zu strukturieren und dennoch mit notwendigen Prozessen, Anweisungen sowie Rahmenpunkten zu stützen}}
	% \printglossary[style=altlist,title=\langglossar]

    % Literaturverzeichnis
	% \clearpage
	\printbibliography
	% \printbibliography[notkeyword=intern, nottype=patent]
	% \printbibliography[heading=subbibliography, keyword=intern, title={Interne Quellen}]
	% \printbibliography[heading=subbibliography, type=patent, title={Patentschriften}]
	% Unterverzeichnisse siehe: https://texwelt.de/fragen/7532/wie-unterteile-ich-meine-biblatex-bibliografie

	% sonstiger Anhang
	% \cleardoublepage
	\clearpage
	\pagenumbering{Alph}
	
	\appendix
	% !TeX root = ../main.tex

% \appendixtoc
\appendix
% \addcontentsline{toc}{chapter}{Appendix}
\ohead[]{\textsc{Appendix}}

% \renewcommand\thechapter{\roman{chapter}}
% \setcounter{chapter}{0}

% \pagebreak
% \includepdf[pages=-,scale=.9,pagecommand={}]{Aufgabenstellung.pdf} 
% PDF um 10% verkleinert einbinden --> Kopf- und Fußzeile  werden so korrekt dargestellt. Die Option `pages' ermöglicht es, eine bestimmte Sequenz von Seiten (z.B. 2-10 oder `-' für alle Seiten) auszuwählen.
% \pagebreak
%\includepdf[pages=-,scale=.8,pagecommand=\section*{A. eventGenerator.py}]{../appendix/eventGenerator.py.pdf}
%\includepdf[pages=-,scale=.8,pagecommand=\section*{B. sendEvents.py}]{../appendix/sendEvents.py.pdf}

%%%%%%%%%%%%%%%%%%%%%%%%%%%%%%%%%%%%%%%%%%
%%%%%%%%%%%%%%%%%%%%%%%%%%%%%%%%%%%%%%%%%%
%%%%%%%%%%%%%%%%%%%%%%%%%%%%%%%%%%%%%%%%%%

\chapter{Anhang A}

Hallo Welt!

\chapter{Anhang B}
\section{Teil 1}

\section{Teil 2}

\end{document}
